\section*{अध्याय \ref{ch:g5:irreversible-changes} --- ऐसे परिवर्तन जो उलटे नहीं किए जा सकते}

\begin{solution}{exo:g5:irreversible-changes:1}
चॉकलेट पिघलाना: उलटा किया जा सकता है (ठंडा करने पर वह फिर जम जाती है)।
तीली जलाना: नहीं। नमक घोलना: उलटा किया जा सकता है (पानी का \omterm{def:g3:separating-mixtures:evaporate}{वाष्पन} करो)।
डबलरोटी सेंकना: नहीं।
\end{solution}

\begin{solution}{exo:g5:irreversible-changes:2}
ऐसा परिवर्तन जिसमें नए पदार्थ प्रकट होते हैं। रसोई में: अंडा पकाना,
केक सेंकना (और भी: डबलरोटी सेंकना, कटे सेब का भूरा होना)।
\end{solution}

\begin{solution}{exo:g5:irreversible-changes:3}
पानी का \omterm{def:g3:separating-mixtures:evaporate}{वाष्पन} करके: नमक प्याली में रह जाता है। \omterm{def:g2:mixing-and-dissolving:dissolve}{घुलना} एक
\omterm{def:g5:irreversible-changes:reversible}{उत्क्रमणीय परिवर्तन} है।
\end{solution}

\begin{solution}{exo:g5:irreversible-changes:4}
नई गंध, और दूध में गुठलियाँ (एक ठोस) प्रकट होना।
\end{solution}

\begin{solution}{exo:g5:irreversible-changes:5}
बाईं ओर के परिवर्तन: पानी में चीनी का \omterm{def:g2:mixing-and-dissolving:dissolve}{घुलना}, बर्फ़ का पिघलना, रेत
और बजरी का मिलाया जाना। दोहरे तीर का अर्थ है कि परिवर्तन दोनों ओर
हो सकता है: उसे उलटा किया जा सकता है।
\end{solution}

\begin{solution}{exo:g5:irreversible-changes:6}
तेल \omterm{def:g4:air-a-mixture-of-gases:air}{हवा} और पानी को चेन के लोहे से दूर रखता है: वह जंग लगने को धीमा
करता है।
\end{solution}

\begin{solution}{exo:g5:irreversible-changes:7}
बिना गरम किए किसी द्रव में गैस के बुलबुले प्रकट होना।
\end{solution}

\begin{solution}{exo:g5:irreversible-changes:8}
उपयोगी: खाना पकाना, डबलरोटी सेंकना, तापन के लिए गैस जलाना। हानिकारक:
जंग लगना, खाना ख़राब होना, लकड़ी का सड़ना।
\end{solution}

\begin{solution}{exo:g5:irreversible-changes:9}
ठंडक उन \omterm{def:g5:irreversible-changes:chemical-change}{रासायनिक परिवर्तनों} को धीमा कर देती है जिनसे दूध खट्टा होता है,
इसलिए वह ज़्यादा दिन टिकता है।
\end{solution}

\begin{solution}{exo:g5:irreversible-changes:10}
नहीं। बुलबुले जलवाष्प हैं, जो ठंडी होने पर फिर पानी बन जाती है: कुछ
नया नहीं बनता, और परिवर्तन को उलटा किया जा सकता है। अकेले बुलबुले यह
सिद्ध नहीं करते कि कोई नया पदार्थ प्रकट हुआ है।
\end{solution}

\begin{solution}{exo:g5:irreversible-changes:11}
\omterm{def:g4:what-a-fire-needs:burning}{जलती} मोमबत्ती ऊष्मा और प्रकाश देती है, और नए पदार्थ बनाती है:
जलवाष्प (उसके ऊपर ठंडा गिलास धुँधला हो जाता है) और कार्बन डाइऑक्साइड।
जो मोम जल जाता है, वह वापस नहीं आता। दूसरी ओर, पिघला हुआ मोम ठंडा होने पर
फिर जम जाता है।
\end{solution}
